% GENERATED FILE. Edit canonical lessons, then run npm run generate:books.
% canonical-source-hash: fnv1a64:d825d6b49a3df5e8
% canonical-lessons: ML-C119-slowly, ML-C119-oncemore, ML-C119-loudly, ML-C119-meaning, ML-C119-wait

\chapter{Slowly, Once More}
\label{ch:ml-slowly-once-more}
\hlchaptermodality{119}

\begin{chapteropening}
\textbf{By the end of this chapter:} \emph{I can ask someone to speak slowly, say it once more, speak louder, give the meaning, and wait.}

Complete the last lesson of chapter 119: I can ask someone to speak slowly, say it once more, speak louder, give the meaning, and wait.
\end{chapteropening}

\section[\texorpdfstring{\ml{പതുക്കെ} (patukke)}{പതുക്കെ (patukke)}]{\ml{പതുക്കെ} (patukke) — slowly}

\label{lesson:ML-C119-slowly}



\begin{quote}
Before the new one: say the Malayalam for mist, then the Malayalam for thunder and lightning.
\end{quote}

\subsection*{You'll want to know: \ml{പതുക്കെ}}
\textbf{\ml{പതുക്കെ}} (\emph{patukke}) — \textquotedblleft{}slowly\textquotedblright{}. \textbf{\ml{പതുക്കെ} \ml{പറയൂ}} (\emph{patukke paṟayū}) — \textquotedblleft{}please say it slowly\textquotedblright{}.

\begin{cousinweb}
\textbf{\ml{പതുക്കെ}} (\emph{patukke}) is native Dravidian, from the verb \textquotedblleft{}to crouch, move stealthily\textquotedblright{}; DEDR 3912 lists the root with Tamil \emph{patuṅku} and Kannada \emph{padugu}.
\end{cousinweb}

\begin{grammarlens}[title={What you've built}]
The most useful request a learner can make.
\end{grammarlens}

\subsection*{Your turn}
\begin{itemize}
\raggedright
  \item \emph{Say it:} \emph{patukke}
  \item \emph{Say it:} \emph{patukke}, once more
  \item \emph{Say it:} say \emph{patukke paṟayū}
  \item \emph{Recall:} say the Malayalam for mist, then the Malayalam for thunder and lightning, then say \emph{patukke} again
\end{itemize}

\begin{tcolorbox}[breakable,colback=teal!4,colframe=teal!35!black,title={Before you move on}]
What does \ml{പതുക്കെ} mean? (\textquotedblleft{}Slowly\textquotedblright{}.) Say it once more.
\end{tcolorbox}
\section[\texorpdfstring{\ml{ഒന്നുകൂടി} (onnukūṭi)}{ഒന്നുകൂടി (onnukūṭi)}]{\ml{ഒന്നുകൂടി} (onnukūṭi) — once more}

\label{lesson:ML-C119-oncemore}



\begin{quote}
Before the new one: say the Malayalam for thunder and lightning, then the Malayalam for slowly.
\end{quote}

\subsection*{You'll want to know: \ml{ഒന്നുകൂടി}}
\textbf{\ml{ഒന്നുകൂടി}} (\emph{onnukūṭi}) — \textquotedblleft{}once more\textquotedblright{}, \textbf{\ml{ഒന്ന്}} (\emph{onnŭ}), \textquotedblleft{}one\textquotedblright{}, with \textbf{\ml{കൂടി}}, \textquotedblleft{}also, more\textquotedblright{}. \textbf{\ml{ഒന്നുകൂടി} \ml{പറയൂ}} — \textquotedblleft{}say it once more\textquotedblright{}.

\begin{cousinweb}
\textbf{\ml{ഒന്നുകൂടി}} (\emph{onnukūṭi}) is native Dravidian: \emph{onnŭ}, \textquotedblleft{}one\textquotedblright{} (DEDR 990), with \emph{kūṭi}, \textquotedblleft{}joined, also\textquotedblright{}, from the verb \textquotedblleft{}to come together\textquotedblright{} (DEDR 1882, with Tamil \emph{kūṭu}).
\end{cousinweb}

\begin{grammarlens}[title={What you've built}]
One, and more.
\end{grammarlens}

\subsection*{Your turn}
\begin{itemize}
\raggedright
  \item \emph{Say it:} \emph{onnukūṭi}
  \item \emph{Say it:} \emph{onnukūṭi}, once more
  \item \emph{Say it:} say \emph{onnukūṭi paṟayū}
  \item \emph{Recall:} say the Malayalam for thunder and lightning, then the Malayalam for slowly, then say \emph{onnukūṭi} again
\end{itemize}

\begin{tcolorbox}[breakable,colback=teal!4,colframe=teal!35!black,title={Before you move on}]
What does \ml{ഒന്നുകൂടി} mean? (\textquotedblleft{}Once more\textquotedblright{}.) Say it once more.
\end{tcolorbox}
\section[\texorpdfstring{\ml{ഉറക്കെ} (uṟakke)}{ഉറക്കെ (uṟakke)}]{\ml{ഉറക്കെ} (uṟakke) — loudly}

\label{lesson:ML-C119-loudly}



\begin{quote}
Before the new one: say the Malayalam for slowly, then the Malayalam for once more.
\end{quote}

\subsection*{You'll want to know: \ml{ഉറക്കെ}}
\textbf{\ml{ഉറക്കെ}} (\emph{uṟakke}) — \textquotedblleft{}loudly, aloud\textquotedblright{}. \textbf{\ml{കുറച്ച്} \ml{ഉറക്കെ}} — \textquotedblleft{}a little louder\textquotedblright{}, with \textbf{\ml{കുറച്ച്}}, \textquotedblleft{}a little\textquotedblright{}.

\begin{cousinweb}
\textbf{\ml{ഉറക്കെ}} (\emph{uṟakke}) is native Dravidian: DEDR 721 lists Malayalam \emph{uṟakkē}, \textquotedblleft{}strongly, firmly, aloud\textquotedblright{}, under the verb \textquotedblleft{}to be firm\textquotedblright{}, with Tamil \emph{uṟai}, \textquotedblleft{}to become firm\textquotedblright{}, and Telugu \emph{oṟapu}, \textquotedblleft{}firmness\textquotedblright{}.
\end{cousinweb}

\begin{grammarlens}[title={What you've built}]
A second way of saying how.
\end{grammarlens}

\subsection*{Your turn}
\begin{itemize}
\raggedright
  \item \emph{Say it:} \emph{uṟakke}
  \item \emph{Say it:} \emph{uṟakke}, once more
  \item \emph{Say it:} say \emph{kuṟaccŭ uṟakke}
  \item \emph{Recall:} say the Malayalam for slowly, then the Malayalam for once more, then say \emph{uṟakke} again
\end{itemize}

\begin{tcolorbox}[breakable,colback=teal!4,colframe=teal!35!black,title={Before you move on}]
What does \ml{ഉറക്കെ} mean? (\textquotedblleft{}Loudly\textquotedblright{}.) Say it once more.
\end{tcolorbox}
\section[\texorpdfstring{\ml{അർത്ഥം} (arthaṁ)}{അർത്ഥം (arthaṁ)}]{\ml{അർത്ഥം} (arthaṁ) — meaning}

\label{lesson:ML-C119-meaning}



\begin{quote}
Before the new one: say the Malayalam for once more, then the Malayalam for loudly.
\end{quote}

\subsection*{You'll want to know: \ml{അർത്ഥം}}
\textbf{\ml{അർത്ഥം}} (\emph{arthaṁ}) — \textquotedblleft{}meaning\textquotedblright{}. \textbf{\ml{ഇതിന്റെ} \ml{അർത്ഥം} \ml{എന്താണ്}?} — \textquotedblleft{}what is the meaning of this?\textquotedblright{}, with \textbf{\ml{എന്താണ്}}, \textquotedblleft{}what is\textquotedblright{}.

\begin{cousinweb}
\textbf{\ml{അർത്ഥം}} (\emph{arthaṁ}) $\leftarrow$ Sanskrit \emph{artha}, \textquotedblleft{}aim, thing, meaning\textquotedblright{}, taken over whole as a tatsama, with Malayalam's noun ending \emph{-ṁ} (CDIAL 638).
\end{cousinweb}

\begin{grammarlens}[title={What you've built}]
The question behind every new word.
\end{grammarlens}

\subsection*{Your turn}
\begin{itemize}
\raggedright
  \item \emph{Say it:} \emph{arthaṁ}
  \item \emph{Say it:} \emph{arthaṁ}, once more
  \item \emph{Say it:} ask \emph{itinṟe arthaṁ entāṇŭ?}
  \item \emph{Recall:} say the Malayalam for once more, then the Malayalam for loudly, then say \emph{arthaṁ} again
\end{itemize}

\begin{tcolorbox}[breakable,colback=teal!4,colframe=teal!35!black,title={Before you move on}]
What does \ml{അർത്ഥം} mean? (\textquotedblleft{}Meaning\textquotedblright{}.) Say it once more.
\end{tcolorbox}
\section[\texorpdfstring{\ml{കാത്തിരിക്കൂ} (kāttirikkū)}{കാത്തിരിക്കൂ (kāttirikkū)}]{\ml{കാത്തിരിക്കൂ} (kāttirikkū) — please wait}

\label{lesson:ML-C119-wait}



\begin{quote}
Before the new one: say the Malayalam for loudly, then the Malayalam for meaning.
\end{quote}

\subsection*{You'll want to know: \ml{കാത്തിരിക്കൂ}}
\textbf{\ml{കാത്തിരിക്കൂ}} (\emph{kāttirikkū}) — \textquotedblleft{}please wait\textquotedblright{}. The ending \textbf{‑\ml{ഊ}} (\emph{-ū}) makes a polite request, as in \textbf{\ml{പറയൂ}}.

\begin{cousinweb}
\textbf{\ml{കാത്തിരിക്കൂ}} (\emph{kāttirikkū}) is native Dravidian, \textquotedblleft{}sit watching\textquotedblright{}: a form of the verb \textquotedblleft{}to guard, watch, wait\textquotedblright{} (DEDR 1416, with Tamil \emph{kā}) joined to the verb \textquotedblleft{}to sit, to be\textquotedblright{} (DEDR 480, with Tamil \emph{iru}), in its request form.
\end{cousinweb}

\begin{grammarlens}[title={What you've built}]
A polite request ending you know.
\end{grammarlens}

\subsection*{Your turn}
\begin{itemize}
\raggedright
  \item \emph{Say it:} \emph{kāttirikkū}
  \item \emph{Say it:} \emph{kāttirikkū}, once more
  \item \emph{Say it:} say \emph{kāttirikkū}
  \item \emph{Recall:} say the Malayalam for loudly, then the Malayalam for meaning, then say \emph{kāttirikkū} again
\end{itemize}

\begin{tcolorbox}[breakable,colback=teal!4,colframe=teal!35!black,title={Before you move on}]
What does \ml{കാത്തിരിക്കൂ} mean? (\textquotedblleft{}Please wait\textquotedblright{}.) Say it once more.
\end{tcolorbox}
